\documentclass[12pt]{exam}
\usepackage[margin=0.1in]{geometry}

\usepackage{amsmath,amssymb}
\usepackage{graphicx}
\usepackage{enumitem}
\usepackage{multicol}
\usepackage{xcolor}
\usepackage{tikz}
\usetikzlibrary{arrows.meta,positioning,shapes.callouts}

\newcommand{\tfq}[3]{%
  \question[#1]\textbf{T/F:} #2%
  % Only show the choices when printing answers/keys
  \ifprintanswers
    \par\begin{oneparchoices}
      #3
    \end{oneparchoices}
  \fi
}



\makeatletter
\newcommand{\TFQ@choices}{} % temp holder

\newenvironment{tfqverb}[2]{%
  \def\TFQ@choices{#2}% store choices for use at end
  \question[#1]\textbf{T/F:}%
}{%
  \ifprintanswers
    \par\begin{oneparchoices}\TFQ@choices\end{oneparchoices}
  \fi
}
\makeatother

\printanswers
% optional styling for the correct choice:
\CorrectChoiceEmphasis{\bfseries}    % make the correct choice bold
% or e.g. \CorrectChoiceEmphasis{\color{red}\bfseries}

\usepackage{fancyvrb}
\newsavebox{\verbbox}

\begin{document}


\begin{center}
\textbf{\Large Test 2 — Database Management Systems}

\textbf{CSC 3300 \quad Spring 2026}
\end{center}

\begin{questions}
%---------------------------------TRUE/FALSE---------------------------------
\section*{Part I — True/False}

\begin{minipage}{\linewidth}
\question[2]
True or False: Consider a view created with the following command.\par\vspace{0.5em}
\begin{verbatim}
create view faculty as
select id, name, dept_name
from instructor
\end{verbatim}
The following two queries return the same data for any instance of the University database.\par\vspace{0.5em}
\begin{verbatim}
select count(*) from instructor
select count(*) from faculty
\end{verbatim}
\begin{checkboxes}
    \choice True
    \CorrectChoice False
\end{checkboxes}
\end{minipage}\vspace{1em}

\begin{minipage}{\linewidth}
\question[2]
True or False: Once a transaction has executed commit work, its effects can no longer be undone by rollback work.\par\vspace{0.5em}
\begin{checkboxes}
    \CorrectChoice True
    \choice False
\end{checkboxes}
\end{minipage}\vspace{1em}

\begin{minipage}{\linewidth}
\question[2]
True or False: If we try to delete a tuple from the \textbf{advisor} relation, the tuple is not deleted. Instead, \textbf{i\_id} of this tuple is set to NULL.\par\vspace{0.5em}
\begin{checkboxes}
    \choice True
    \CorrectChoice False
\end{checkboxes}
\end{minipage}\vspace{1em}

\begin{minipage}{\linewidth}
\question[2]
True or False: Definition of the \textit{prereq} table could be changed from: \par\vspace{0.5em}
\begin{verbatim}
create table prereq
(course_id        varchar(8), 
prereq_id        varchar(8),
primary key (course_id, prereq_id),
foreign key (course_id) references course(course_id) on delete cascade,
foreign key (prereq_id) references course(course_id)
);
\end{verbatim}
to\par\vspace{0.5em}
\begin{verbatim}
create table prereq
(course_id        varchar(8), 
prereq_id        varchar(8),
primary key (course_id, prereq_id),
foreign key (course_id) references course(course_id) on delete cascade,
foreign key (prereq_id) references course(course_id) on delete set null
);
\end{verbatim}
\begin{checkboxes}
    \choice True
    \CorrectChoice False
\end{checkboxes}
\end{minipage}\vspace{1em}

\begin{minipage}{\linewidth}
\question[2]
True or False: MySQL Connector/J implements the JDBC API.\par\vspace{0.5em}
\begin{checkboxes}
    \CorrectChoice True
    \choice False
\end{checkboxes}
\end{minipage}\vspace{1em}

\begin{minipage}{\linewidth}
\question[2]
True or False: Materialized views physically store the result set of a query, and their data can be updated periodically.\par\vspace{0.5em}
\begin{checkboxes}
    \CorrectChoice True
    \choice False
\end{checkboxes}
\end{minipage}\vspace{1em}

\begin{minipage}{\linewidth}
\question[2]
True or False: A primary key constraint automatically creates an index.\par\vspace{0.5em}
\begin{checkboxes}
    \CorrectChoice True
    \choice False
\end{checkboxes}
\end{minipage}\vspace{1em}

\begin{minipage}{\linewidth}
\question[2]
True or False: Deleting a tuple from the \textbf{instructor} relation can result in setting value(s) of some attribute(s) 
of some tuples in the \textbf{advisor} relation to NULL.\par\vspace{0.5em}
\begin{checkboxes}
    \CorrectChoice True
    \choice False
\end{checkboxes}
\end{minipage}\vspace{1em}


%---------------------------------MC---------------------------------
\section*{Part II — Multiple choice (choose one)}

\begin{minipage}{\linewidth}
\question[3]
What is the maximal number of tables from which tuple(s) can be automatically deleted with the deletion of a \textbf{student}?\par\vspace{0.5em}
\begin{checkboxes}
    \choice 1
    \choice 0
    \CorrectChoice 2
    \choice 4
\end{checkboxes}
\end{minipage}\vspace{1em}

\begin{minipage}{\linewidth}
\question[3]
Which code snippet is vulnerable to SQL injection attacks?\par\vspace{0.5em}
\begin{checkboxes} 
    \CorrectChoice \begin{verbatim}
String id = "225";
String name = "Green";
String dept_name = "Comp. Sci.";
BigDecimal salary = new BigDecimal(98000);
Statement stmt = conn.createStatement();
stmt.executeUpdate("insert into instructor values 
('" + id + "', '" + name + "', '" + dept_name + "', " + salary + ")");
    \end{verbatim} %-- String name = "Green', 'Comp. Sci.', 3) -- ";
    \choice \begin{verbatim}
String id = "225";
String name = "Green";
String dept_name = "Comp. Sci.";
BigDecimal salary = new BigDecimal(98000);
PreparedStatement pstmt = conn.prepareStatement
("insert into instructor values (?, ?, ?, ?)");
pstmt.setString(1, id);
pstmt.setString(2, name);
pstmt.setString(3, dept_name);
pstmt.setBigDecimal(4, salary);
pstmt.executeUpdate();
    \end{verbatim}
\end{checkboxes}
\end{minipage}\vspace{1em}

%---------------------------------MS---------------------------------
\section*{Part III — Multiple select (select all that apply)}

\begin{minipage}{\linewidth}
\question[5]
Consider the following view, named myView, defined in the University database.\par\vspace{0.5em}
\begin{verbatim}
create view myView(dept_name, course_id, amount) as
       select dept_name, course_id, count(*)
       from course natural join section
       group by dept_name, course_id;
\end{verbatim}
Which of the given SQL statements are allowed on the view?\par\vspace{0.5em}
\begin{checkboxes}
    \CorrectChoice select dept\_name, course\_id from myView
    \CorrectChoice select * from course where course\_id in (select course\_id from myView)
    \CorrectChoice select course\_id, amount from myView where amount \texttt{=} (select min(amount) from myView)
    \choice create view nextView(course\_id) as select course\_id, title from myView where amount \texttt{>=} 2
    \choice insert into myView values ('Physics', 'PHY-101', 1)
\end{checkboxes}
\end{minipage}\vspace{1em}

\begin{minipage}{\linewidth}
\question[4]
Which statements are true regarding the following procedure?\par\vspace{0.5em}

\begin{verbatim}
create procedure instructors_no_proc
(in dept_name varchar(20), out d_count integer)
begin
select count(*) into d_count 
from instructor where instructor.dept_name = dept_name;
end
\end{verbatim}

{\footnotesize Note: To create a procedure in MySQL Workbench, wrap the procedure definition with

\texttt{delimiter //} 

before it, and 

\texttt{// delimiter ;}

after it.}

\par\vspace{0.5em}

\begin{checkboxes}
    \CorrectChoice The following command correctly calls the procedure:
    \begin{verbatim}
call instructors_no_proc('Comp. Sci.', @a);
    \end{verbatim}

    \choice The following command correctly calls the procedure:
    \begin{verbatim}
select instructors_no_proc('Comp. Sci.', @a);
    \end{verbatim}

    \choice The following command correctly calls the procedure:
    \begin{verbatim}
call instructors_no_proc('Comp. Sci.', 7);
    \end{verbatim}

    \choice The execution of instructors\_no\_proc procedure can result in a modification to the database data.
\end{checkboxes}
\end{minipage}\vspace{1em}

\end{questions}
\end{document}